\scanpage{93}
% Typescript inverted in the scan.
\mspage{\handwritten{42}}
d'où un homomorphisme $\mathrm{f^{\unsure{*|x}}(G)\mathrel{\relbar\joinrel\handwritten{>}} F'}$ et ce dernier est evidemment
compatible avec les données de descente. Il reste à examiner si cet
homomorphisme est un \uline{isomorphisme.}
\figure{\begin{tikzpicture}[x=1cm,y=1cm]
\draw[line width=.35pt] (0,-.23) .. controls (-.09,-.18) and (-.02,.02) .. (.12,.05) .. controls (.6,.08) and (.79,.05) .. (1.2,.07) .. controls (1.65,.02) and (1.72,.13) .. (2.22,.14) -- (4.24,.18) .. controls (4.53,.18) and (4.59,.24) .. (4.56,.40) .. controls (4.56,.46) and (4.53,.49) .. (4.48,.49);
\end{tikzpicture}}
Noter que $\mathrm{G}$ est aussi définissable par l'exactitude de
\diagram{\begin{tikzcd}[ampersand replacement=\&,column sep=large]
\mathrm{G}\arrow[r,dash,"\handwritten{>}"{pos=1,anchor=east,inner sep=0pt}] \& \mathrm{f_*(F')}\arrow[r,dash,shift left=.6ex,"\handwritten{>}"{pos=1,anchor=east,inner sep=0pt}]\arrow[r,dash,shift right=.6ex,"\handwritten{>}"{pos=1,anchor=east,inner sep=0pt}] \& \mathrm{g_*(F'')}\text{,}
\end{tikzcd}}
donc pour tout changement de base $\mathrm{Y\overset{h}{\mathrel{\relbar\joinrel\handwritten{>}}}X}$ \uline{qui commute à la formation
de} $\mathrm{f_*(F')}$ \uline{et} $\mathrm{g_*(F'')}$\handwritten{,} $\mathrm{h^*(G)=G_Y}$ s'identifie au faisceau ``descendu''
de $\mathrm{F'_{Y'}}$ par $\mathrm{Y' \mathrel{\relbar\joinrel\handwritten{>}}Y}$, et bien entendu l'homomorphisme
$\mathrm{f_Y^*(G_Y)\mathrel{\relbar\joinrel\handwritten{>}}(F_Y)'}$ est celui déduit de $\mathrm{f^*(G)\mathrel{\relbar\joinrel\handwritten{>}}F'}$ par changement
de base. Or supposons que \uline{le changement de base} $\mathrm{f}$ ; $\mathrm{Y=X'\mathrel{\relbar\joinrel\handwritten{>}}X}$
\uline{commute à la formation de} $\mathrm{f_*(F')}$ \uline{et} $\mathrm{g_*(F'')}$\handwritten{,} et notons que comme
$\mathrm{Y'=X' \handwritten{\times}_X X'}$ a une section sur $\mathrm{Y'=X'}$, donc $\mathrm{Y'\mathrel{\relbar\joinrel\handwritten{>}}Y}$ est un
morphisme de descente \uline{effective} pour toute catégorie fibrée, en particulier
pour $\mathrm{\handwritten{\unsure{\mathfrak F|\mathcal F}}}$, il s'ensuit que $\mathrm{f_{\deleted{\unsure{X}}\handwritten{Y}}^*(G_{\deleted{\unsure{X}}\handwritten{Y}})\mathrel{\relbar\joinrel\handwritten{>}}(F_{\deleted{\unsure{X}}\handwritten{Y}})'}$ est un isomorphisme,
donc $\mathrm{f^*(G)\mathrel{\relbar\joinrel\handwritten{>}}F'}$ dévient un isomorphisme après changement de
base $\mathrm{Y'=X'\quad {}_X X'\mathrel{\relbar\joinrel\handwritten{>}}X'}$. Comme ce dernier à une section,
$\mathrm{f^*(G)\mathrel{\relbar\joinrel\handwritten{>}}F'}$ est un isomorphisme, donc la donnée de descente envisagée
sur $\mathrm{F'}$ est effective.

9.4.2. Ceci prouve le théorème 9.4. dans le cas a), grâce à 5.6.
Le cas b) résulte également du fait que si $\mathrm{f}$ est propre, alors $\mathrm{f_*}$
commute à \uline{tout} changement de base $\mathrm{Y' \mathrel{\relbar\joinrel\handwritten{>}} Y}$, (qui sera prouvé\unsure{'} dans un
exposé ultérieur).

9.4.3. Dans le cas c), la question étant locale sur $\mathrm{\unsure{Y|Y'}}$ on peut
% archive sheet number: 92

